\documentclass[paper=a4paper,fontsize=10pt,jafontsize=9pt,titlepage]{jlreq}
%=================================================================
% LuaLaTeX preamble
% (%=================================================================
% LuaLaTeX preamble
% (%=================================================================
% LuaLaTeX preamble
% (\input{preamble}でtemplate.texに読み込まれることを想定)
%
% 用途:
%   日本語・英語・古典ギリシア語混在文書
%   (哲学・神学・古典学向け)
%
% 2026-06
%
% 重要:
%   - Unicode Greek を使用可能
%   - Babel/LGR の古い \textgreek{...} 入力を維持
%   - hyperref による PDF しおり対応
%=================================================================

%=================================================================
% 日本語処理
%=================================================================

\usepackage{luatexja-fontspec}
\usepackage{luatexja-ruby}

% 古典ギリシア語(Unicode)を和文扱いしない。
%
% これが無いと
%
%   λόγος
%   ἀλήθεια
%
% などが HaranoAjiMincho に送られ、
%
%   Missing character: There is no ό ...
%
% が発生する。
%
\ltjsetparameter{jacharrange={-2,-3}}

%=================================================================
% フォント
%=================================================================

% 欧文・ギリシア語
\setmainfont{Libertinus Serif}

% 和文
\setmainjfont{HaranoAjiMincho}

%=================================================================
% 古典ギリシア語
%=================================================================

% Babel の古い LGR 入力
%
%   \textgreek{l'ogos}
%   \textgreek{>'anjrwpoc}
%
% を維持するために必要。
%
% Unicode Greek
%
%   λόγος
%   \foreignlanguage{greek}{πάντες ἄνθρωποι}
%
% と共存可能。
%
%  しかしitalic, boldなどが出なくなるので、fontencは外す。その結果、LGRは不可
% \usepackage[LGR,T1]{fontenc}
\usepackage[polutonikogreek,english,japanese]{babel}

%=================================================================
% レイアウト
%=================================================================

\usepackage{paracol}
\footnotelayout{m}

\usepackage{longtable}

%=================================================================
% 色
%=================================================================

\usepackage{xcolor}

%=================================================================
% PDF ハイパーリンク
%=================================================================

\usepackage[
unicode,
bookmarks=true,
hidelinks
]{hyperref}

% hyperref の後に読むこと。
% 先に読むと Option clash が発生する。
%
\usepackage{bookmark}

%=================================================================
% ヘッダ・フッタ
%=================================================================

\usepackage{fancyhdr}
\pagestyle{fancy}
%=================================================================
% 注意事項
%=================================================================
%
% 見出し中のギリシア語:
%
%   \section{λόγος}
%
% は空白になる場合がある。
%
% その場合:
%
%   \section{{\rmfamily λόγος}}
%
% と書く。
%
%
% Babel ギリシア語:
%
%   \textgreek{l'ogos}
%
% Unicode ギリシア語:
%
%   \foreignlanguage{greek}{λόγος}
%
% の両方が利用可能。
%
%=================================================================
でtemplate.texに読み込まれることを想定)
%
% 用途:
%   日本語・英語・古典ギリシア語混在文書
%   (哲学・神学・古典学向け)
%
% 2026-06
%
% 重要:
%   - Unicode Greek を使用可能
%   - Babel/LGR の古い \textgreek{...} 入力を維持
%   - hyperref による PDF しおり対応
%=================================================================

%=================================================================
% 日本語処理
%=================================================================

\usepackage{luatexja-fontspec}
\usepackage{luatexja-ruby}

% 古典ギリシア語(Unicode)を和文扱いしない。
%
% これが無いと
%
%   λόγος
%   ἀλήθεια
%
% などが HaranoAjiMincho に送られ、
%
%   Missing character: There is no ό ...
%
% が発生する。
%
\ltjsetparameter{jacharrange={-2,-3}}

%=================================================================
% フォント
%=================================================================

% 欧文・ギリシア語
\setmainfont{Libertinus Serif}

% 和文
\setmainjfont{HaranoAjiMincho}

%=================================================================
% 古典ギリシア語
%=================================================================

% Babel の古い LGR 入力
%
%   \textgreek{l'ogos}
%   \textgreek{>'anjrwpoc}
%
% を維持するために必要。
%
% Unicode Greek
%
%   λόγος
%   \foreignlanguage{greek}{πάντες ἄνθρωποι}
%
% と共存可能。
%
%  しかしitalic, boldなどが出なくなるので、fontencは外す。その結果、LGRは不可
% \usepackage[LGR,T1]{fontenc}
\usepackage[polutonikogreek,english,japanese]{babel}

%=================================================================
% レイアウト
%=================================================================

\usepackage{paracol}
\footnotelayout{m}

\usepackage{longtable}

%=================================================================
% 色
%=================================================================

\usepackage{xcolor}

%=================================================================
% PDF ハイパーリンク
%=================================================================

\usepackage[
unicode,
bookmarks=true,
hidelinks
]{hyperref}

% hyperref の後に読むこと。
% 先に読むと Option clash が発生する。
%
\usepackage{bookmark}

%=================================================================
% ヘッダ・フッタ
%=================================================================

\usepackage{fancyhdr}
\pagestyle{fancy}
%=================================================================
% 注意事項
%=================================================================
%
% 見出し中のギリシア語:
%
%   \section{λόγος}
%
% は空白になる場合がある。
%
% その場合:
%
%   \section{{\rmfamily λόγος}}
%
% と書く。
%
%
% Babel ギリシア語:
%
%   \textgreek{l'ogos}
%
% Unicode ギリシア語:
%
%   \foreignlanguage{greek}{λόγος}
%
% の両方が利用可能。
%
%=================================================================
でtemplate.texに読み込まれることを想定)
%
% 用途:
%   日本語・英語・古典ギリシア語混在文書
%   (哲学・神学・古典学向け)
%
% 2026-06
%
% 重要:
%   - Unicode Greek を使用可能
%   - Babel/LGR の古い \textgreek{...} 入力を維持
%   - hyperref による PDF しおり対応
%=================================================================

%=================================================================
% 日本語処理
%=================================================================

\usepackage{luatexja-fontspec}
\usepackage{luatexja-ruby}

% 古典ギリシア語(Unicode)を和文扱いしない。
%
% これが無いと
%
%   λόγος
%   ἀλήθεια
%
% などが HaranoAjiMincho に送られ、
%
%   Missing character: There is no ό ...
%
% が発生する。
%
\ltjsetparameter{jacharrange={-2,-3}}

%=================================================================
% フォント
%=================================================================

% 欧文・ギリシア語
\setmainfont{Libertinus Serif}

% 和文
\setmainjfont{HaranoAjiMincho}

%=================================================================
% 古典ギリシア語
%=================================================================

% Babel の古い LGR 入力
%
%   \textgreek{l'ogos}
%   \textgreek{>'anjrwpoc}
%
% を維持するために必要。
%
% Unicode Greek
%
%   λόγος
%   \foreignlanguage{greek}{πάντες ἄνθρωποι}
%
% と共存可能。
%
%  しかしitalic, boldなどが出なくなるので、fontencは外す。その結果、LGRは不可
% \usepackage[LGR,T1]{fontenc}
\usepackage[polutonikogreek,english,japanese]{babel}

%=================================================================
% レイアウト
%=================================================================

\usepackage{paracol}
\footnotelayout{m}

\usepackage{longtable}

%=================================================================
% 色
%=================================================================

\usepackage{xcolor}

%=================================================================
% PDF ハイパーリンク
%=================================================================

\usepackage[
unicode,
bookmarks=true,
hidelinks
]{hyperref}

% hyperref の後に読むこと。
% 先に読むと Option clash が発生する。
%
\usepackage{bookmark}

%=================================================================
% ヘッダ・フッタ
%=================================================================

\usepackage{fancyhdr}
\pagestyle{fancy}
%=================================================================
% 注意事項
%=================================================================
%
% 見出し中のギリシア語:
%
%   \section{λόγος}
%
% は空白になる場合がある。
%
% その場合:
%
%   \section{{\rmfamily λόγος}}
%
% と書く。
%
%
% Babel ギリシア語:
%
%   \textgreek{l'ogos}
%
% Unicode ギリシア語:
%
%   \foreignlanguage{greek}{λόγος}
%
% の両方が利用可能。
%
%=================================================================

\lhead{{\itshape Summa Theologiae} I, q.84}
%--------------------


\bibliographystyle{jplain}


\title{{\bfseries Prima Pars}\\{\Huge Summae Theologiae}\\Sancti Thomae
Aquinatis\\{\sffamily QUAESTIO OCTOGESIMAQUARTA}\\{\bfseries QUOMODO ANIMA OCNIUNCTA INTELLIGAT COPORALIA, QUAE SUNT INFRA IPSAM}}
\author{Japanese translation\\by Yoshinori {\scshape Ueeda}}
\date{Last modified \today}

%%%% コピペ用
%\rhead{a.~}
%\begin{center}
% {\Large {\bfseries }}\\
% {\large }\\
% {\footnotesize }\\
% {\Large \\}
%\end{center}
%
%\begin{longtable}{p{21em}p{21em}}
%
%&
%
% 
%
%\\
%\end{longtable}
%\newpage

\begin{document}

\maketitle

\begin{center}
 {\Large 第八十四問\\結合した魂は\\自らの下位にある物体的なものを\\どのように知性認識するか}
\end{center}

\begin{longtable}{p{21em}p{21em}}
Consequenter considerandum est de actibus animae, quantum ad potentias
intellectivas et appetitivas, aliae enim animae potentiae non pertinent directe
ad considerationem theologi. Actus autem appetitivae partis ad considerationem
moralis scientiae pertinent, et ideo in secunda parte huius operis de eis
tractabitur, in qua considerandum erit de morali materia. Nunc autem de actibus
 intellectivae partis agetur.

 &

 続いて、魂の諸々の作用について考察されるべきである。そしてそれは知性的な能力と
 欲求的な能力にかんしてである。なぜなら、これ以外の魂の能力は、直接的に神学者の
 考察に属さないからである。さて、欲求的部分の作用は道徳的学に属する。ゆえに、こ
 の著作の第二部でそれらについて論じられる。そこにおいては、道徳的な材料が考察さ
 れなければならないだろう。他方で、今は、知性的な部分の作用について論じられる。

 \\[1em]

 In consideratione vero actuum, hoc modo procedemus,
primo namque considerandum est quomodo intelligit anima corpori coniuncta;
 secundo, quomodo intelligit a corpore separata.

 &

 さて、諸作用の考察において、私たちは以下のように進む。第一に、身体に結合された
 魂が、どのようにして知性認識するか。第二に、身体から分離した魂が、どのように知
 性認識するか。

 \\[1em]

 Prima autem consideratio erit tripartita, primo namque considerabitur quomodo
anima intelligit corporalia, quae sunt infra ipsam; secundo, quomodo intelligit
seipsam, et ea quae in ipsa sunt; tertio, quomodo intelligit substantias
 immateriales, quae sunt supra ipsam.

 &

 第一の考察は三つの部分を持つ。第一に、どのようにして魂が自らの下位にある物体的
 なものを知性認識するか。第二に、どのようにして自己自身を、そして自己の中にある
 者どもを知性認識するか、第三に、どのようにして自らの上位にある非質料的諸実態を
 知性認識するか。

\\[1em]

 Circa cognitionem vero corporalium, tria consideranda occurrunt, primo quidem
per quid ea cognoscit; secundo, quomodo et quo ordine; tertio, quid in eis
cognoscit. Circa primum quaeruntur octo.

 &

 物体的なものの認識を巡って、三つの考察されるべきことが生じる。第一に、何によっ
 てそれらを認識するか。第二に、どのようにして、どのような順番で認識するか。第三
 に、何をそれらにおいて認識するか。第一をめぐって、八つのことが問われる。

 \\[1em]
 \begin{enumerate}
  \item utrum anima cognoscat corpora per intellectum.
  \item utrum intelligat ea per essentiam suam, vel per aliquas species.
  \item si per aliquas species, utrum species omnium intelligibilium sint ei naturaliter innatae.
  \item utrum effluant in ipsam ab aliquibus formis immaterialibus separatis.
  \item utrum anima nostra omnia quae intelligit, videat in rationibus aeternis.
  \item utrum cognitionem intelligibilem acquirat a sensu.
  \item utrum intellectus possit actu intelligere per species intelligibiles quas penes se habet, non convertendo se ad phantasmata.
  \item utrum iudicium intellectus impediatur per impedimentum sensitivarum virtutum.
 \end{enumerate}

 &

 \begin{enumerate}
  \item 魂は物体を知性によって認識するか。
  \item 自らの本質によって認識するか、それとも何らかの形象によってか。
  \item もし何らかの形象によってであるならば、すべての可知的なものの形象は、魂に本性的に植え付けられているか。
  \item (形象は)何らかの離在する非質料的形相から魂へ流入するか。
  \item 私たちの魂は、知性認識するすべてのものを、永遠の理念において見るか。
  \item (魂は)可知的な認識を感覚から獲得するか。
  \item 知性は、表象像へ向き直ることなしに、自らのもとにある可知的形象によって、現実に知性認識することができるか。
  \item 知性の判断は、感覚能力の妨害によって妨害されるか。
 \end{enumerate}
\end{longtable}

\rhead{a.1}
\begin{center}
{\Large {\bfseries ARTICULUS PRIMUS}}\\
{\large UTRUM ANIMA COGNOSCAT CORPORA PER INTELLECTUM}\\
{\footnotesize {\itshape De Verit}., q.10, a.1}\\
{\Large 第一項\\魂は物体を知性によって認識するか}
\end{center}

\begin{longtable}{p{21em}p{21em}}
1. Ad primum sic proceditur. Videtur quod anima non cognoscat corpora per
intellectum. Dicit enim Augustinus, in II Soliloq. quod corpora intellectu
comprehendi non possunt; nec aliquod corporeum nisi sensibus videri
potest. Dicit etiam, XII super Gen. ad Litt., quod visio intellectualis est
eorum quae sunt per essentiam suam in anima. Huiusmodi autem non sunt
corpora. Ergo anima per intellectum corpora cognoscere non potest.
 
 &

第一項の問題へ、議論は以下のように進められる。魂は知性によって物体を認識しないと
思われる。理由は以下の通り。アウグスティヌスは『ソリロクィア』第2巻で「物体、ま
たどんな物体的なものも、感覚によって見られうるのでない限り、知性によって把握され
得ない」と言う。さらに、『創世記逐語注解』第12巻で、知性的な直視は、自らの本質に
よって魂の中にあるものどもに属すると言う。しかし物体はそのようなものではない。ゆ
えに、魂は、知性によって物体を認識することができない。
 
\\[1em] 

2. Praeterea, sicut se habet sensus ad intelligibilia, ita se habet intellectus ad
 sensibilia. Sed anima per sensum nullo modo potest cognoscere spiritualia, quae
 sunt intelligibilia. Ergo nullo modo per intellectum potest cognoscere corpora,
 quae sunt sensibilia.

 &

 さらに、感覚が可知的なものに関係するように、知性は可感的なものに関係する。しか
 るに、魂は感覚によって、霊的なもの、つまり可知的なものを、どうやっても認識でき
 ない。ゆえに、知性によっては、物体的なもの、つまり可感的なものを、どうしても認
 識できない。
 
 \\[1em]

 3. Praeterea, intellectus est necessariorum et semper eodem modo se habentium. Sed
 corpora omnia sunt mobilia, et non eodem modo se habentia. Anima ergo per
 intellectum corpora cognoscere non potest.

 &

 さらに、知性は、必然的なもの、常に同じしかたであるものに関わる。しかるにすべて
 の物体は動きうるし、同じしかたであるわけでもない。ゆえに魂は知性によって物体を
 認識することができない。

 \\[1em]

 Sed contra est quod scientia est in intellectu. Si ergo intellectus non
 cognoscit corpora, sequitur quod nulla scientia sit de corporibus. Et sic
 peribit scientia naturalis, quae est de corpore mobili.

 &

 しかし反対に、知性の中に学がある。ゆえに、もし知性が物体を認識しないならば、物
 体についての学がないことが帰結する。かくして、動きうる物体についてある自然的な
 学が消えることになる。
 
 \\[1em]

 Respondeo dicendum, ad evidentiam huius quaestionis, quod primi philosophi qui
 de naturis rerum inquisiverunt, putaverunt nihil esse in mundo praeter
 corpus. Et quia videbant omnia corpora mobilia esse, et putabant ea in continuo
 fluxu esse, aestimaverunt quod nulla certitudo de rerum veritate haberi posset
 a nobis.

 &

解答する。この問題を明らかにするためには、以下のように言われるべきである。諸事物
の本性について探究した最初の哲学者たちは、世界の中に物体以外のものはないと考えた。
そしてすべての物体は動的であることを見、それらが連続的な流動の中にあると考えたの
で、諸事物の真理についてどんな確実性も、私たちによって持たれ得ないと考えた。
 
\\[1em]

 Quod enim est in continuo fluxu, per certitudinem apprehendi non potest, quia
 prius labitur\footnote{l\={a}bor, lapsus. to slide, slip away from} quam mente
 diiudicetur, sicut Heraclitus dixit quod non est possibile aquam fluvii
 currentis bis tangere, ut recitat philosophus in IV Metaphys.

 &

というのも、連続的な流動の中にあるものは、確実性によって捉えられえないからである。
なぜなら、精神によって判断されるより先に消え失せるからである。ちょうど、哲学者が
『形而上学』第4巻で引くように、ヘラクレイトスが「流れる川の水に二度触れることは
できない」と言ったように。
 
\\[1em]

 His autem superveniens Plato, ut posset salvare certam cognitionem
 veritatis a nobis per intellectum haberi, posuit praeter ista corporalia aliud
 genus entium a materia et motu separatum, quod nominabat species sive ideas,
 per quarum participationem unumquodque istorum singularium et sensibilium
 dicitur vel homo vel equus vel aliquid huiusmodi.

 &

 しかしこの人たちの後にプラトンが来て、真理のある種の認識が私たちによって知性
 を通して持たれるということを救えるように、この物体的なもの以外に、質料と運動か
 ら分離された、存在の別の類を措定し、彼はそれを形象ないしイデアと呼んでいた。そ
 れらを分有することによって、これらの個物や可感的なものの各々は、人間、馬、ある
 いは何かそのようなものと言われる。
 
\\[1em]

 Sic ergo dicebat scientias et definitiones et quidquid ad actum intellectus
 pertinet, non referri ad ista corpora sensibilia, sed ad illa immaterialia et
 separata; ut sic anima non intelligat ista corporalia, sed intelligat horum
 corporalium species separatas.

 &

 それゆえ、このようにして、学知や定義、そして知性の作用に属するものは何であれ、
 これらの可感的物体ではなく、かの非質料的で分離したものに関係づけられると言った。
 また、このようにして、魂はこの物体的なものを知性認識せず、これらの物体的なもの
 の、分離した形象を知性認識すると述べた。
 
\\[1em]


 Sed hoc dupliciter apparet falsum. Primo quidem quia, cum illae species sint
 immateriales et immobiles, excluderetur a scientiis cognitio motus et materiae
 (quod est proprium scientiae naturalis) et demonstratio per causas moventes et
 materiales.

 &

 しかしこれは二通りに偽である。第一に、彼の形象は非質料的で不動なので、(もしプラ
 トンが正しいなら)学から運動と質料の認識が排除されたであろう(が、これは自然的諸
 学に固有のことである)し、動因と質料因による証明も排除されたであろう。

\\[1em]

 Secundo autem, quia derisibile videtur ut, dum rerum
 quae nobis manifestae sunt notitiam quaerimus, alia entia in medium afferamus,
 quae non possunt esse earum substantiae, cum ab eis differant secundum esse, et
 sic, illis substantiis separatis cognitis, non propter hoc de istis
 sensibilibus iudicare possemus.

 &

 さらに第二に、以下のようなことはおかしなことに思われるからである。つまり、私た
 ちは、私たちに明らかである諸事物についての知を探究しているのに、それらの事物の
 実体に属しえない、他の存在者を中間に持ち込むということは。なぜかというと、それ
 らの存在者は、私たちの周りにある諸事物とは存在において異なり、そのようにして、
 分離したそれらの実体が認識されても、そのために、私たちが、これらの可感的なもの
 どもについて判断することはできないからである。

\\[1em]


 Videtur autem in hoc Plato deviasse a veritate, quia, cum aestimaret omnem
 cognitionem per modum alicuius similitudinis esse, credidit quod forma cogniti
 ex necessitate sit in cognoscente eo modo quo est in cognito.

 &

 さらに、この点においてプラトンは真理から逸脱したと思われる。すなわち彼は、すべ
 ての認識を、何らかの類似のしかたによってあると考えたので、認識されたものの形相
 が、必然的に、認識されたものにおいてあるのと同じしかたで、認識する者の中にある
 と信じた。
 
\\[1em]

 Consideravit autem quod forma rei intellectae est in intellectu
 universaliter et immaterialiter et immobiliter, quod ex ipsa operatione
 intellectus apparet, qui intelligit universaliter et per modum necessitatis
 cuiusdam; modus enim actionis est secundum modum formae agentis. Et ideo
 existimavit quod oporteret res intellectas hoc modo in seipsis subsistere,
 scilicet immaterialiter et immobiliter.


 &

 さらに彼は、知性認識された事物の形相が、知性の中に普遍的で非質料的で不動のしか
 たで存在することを考察し、このことは、普遍的に、そして何らかの必然性のあり方に
 よって知性認識する知性の働き自体から明らかなのだが、つまり作用のあり方は、作用
 者の形相のあり方に従うからだが、それゆえ、知性認識された事物は、それ自体におい
 て、このようなしかたで、すなわち非質料的で不動のしかたで自存するのでなければな
 らないと考えた。
 
\\[1em]

 Hoc autem necessarium non est. Quia
 etiam in ipsis sensibilibus videmus quod forma alio modo est in uno sensibilium
 quam in altero, puta cum in uno est albedo intensior, in alio remissior, et in
 uno est albedo cum dulcedine, in alio sine dulcedine.

 &

 しかしこのことは必然的ではない。なぜなら、可感的なものども自体において、私たち
 は、形相が、可感的なものどものあるものにおいて、他のものにおけるのとは違うしか
 たで存在するのを見る。たとえば、あるものにおいては白性がより強く、他のものにお
 いてはより弱くあり、あるものにおいては、白性が甘さと共にあり、他のものにおいて
 は甘さなしにある、という場合のように。

\\[1em]


 Et per hunc etiam modum forma sensibilis alio modo est in re quae est extra
 animam, et alio modo in sensu, qui suscipit formas sensibilium absque materia,
 sicut colorem auri sine auro.

 &

 そしてさらに、このしかたで、可感的形相は、魂の外にある事物においてと、感覚にお
 いてとでは別のしかたである。感覚は、可感的なものの形相を質料なしに保持する。た
 とえば、金の色を、金なしに保持するように。
 

\\[1em]


 Et similiter intellectus species, corporum, quae sunt materiales et
 mobiles, recipit immaterialiter et immobiliter, secundum modum suum, nam
 receptum est in recipiente per modum recipientis. Dicendum est ergo quod anima
 per intellectum cognoscit corpora cognitione immateriali, universali et
 necessaria.

 &

 同様に、知性は、質料的で動的な物体の形象を、自分のあり方に従って、非質料的で不
 動のしかたで受け取る。すなわち、受け取られるものは、受け取るものの中に、受け取
 るもののあり方によって存在する。ゆえに、魂は知性によって、物体を非質料的で普遍的
 で必然的な認識によって認識する、と言われるべきである。

 
\\[1em]

 Ad primum ergo dicendum quod verbum Augustini est intelligendum quantum ad ea
 quibus intellectus cognoscit, non autem quantum ad ea quae cognoscit. Cognoscit
 enim corpora intelligendo, sed non per corpora, neque per similitudines
 materiales et corporeas; sed per species immateriales et intelligibiles, quae
 per sui essentiam in anima esse possunt.

 &

 第一異論に対しては、それゆえ、以下のように言われるべきである。アウグスティヌス
 の言葉は、\kenten{それを}認識するところの\kenten{それ}に関してではなく、
 \kenten{それによって}知性が認識するところのものどもの\kenten{それ}に関して理解
 されるべきである。なぜなら、知性は物体を知性認識することによって認識するが、し
 かし物体によってではなく、また、質料的で物体的な類似によってでもなく、むしろ、
 それ自らの本質によって魂の中に存在することが可能である、非質料的で可知的な形象
 によって認識するからである。
 
 \\[1em]

 Ad secundum dicendum quod, sicut Augustinus dicit XXII de Civit. Dei, non est
 dicendum quod, sicut sensus cognoscit sola corporalia, ita intellectus
 cognoscit sola spiritualia, quia sequeretur quod Deus et Angeli corporalia non
 cognoscerent. Huius autem diversitatis ratio est, quia inferior virtus non se
 extendit ad ea quae sunt superioris virtutis; sed virtus superior ea quae sunt
 inferioris virtutis, excellentiori modo operatur.

 &

 第二異論に対しては、以下のように言われるべきである。アウグスティヌスが『神の国』
 第22巻で言うように、感覚がただ物体的なものだけを認識するように、知性もただ霊的
 なものだけをにんしきする、と言われるべきではない。なぜなら、もしそうならば、神
 と天使たちは、物体的なものを認識しなかっただろうから。また、この違いの理由は、
 下位のちからは、上位のちからに属するものにまで及ばないが、上位のちからは、下位
 のちからに属するものに、より卓越したしかたで働くからである。

 \\[1em]

 Ad tertium dicendum quod omnis motus supponit aliquid immobile, cum enim
 transmutatio fit secundum qualitatem, remanet substantia immobilis; et cum
 transmutatur forma substantialis, remanet materia immobilis. Rerum etiam
 mutabilium sunt immobiles habitudines, sicut Socrates etsi non semper sedeat,
 tamen immobiliter est verum quod, quandocumque sedet, in uno loco manet. Et
 propter hoc nihil prohibet de rebus mobilibus immobilem scientiam habere.

 &

 第三異論に対しては、以下のように言われるべきである。運動はすべてなんらかの不同
 なものをもとにする。なぜなら、性質に即して変化が生じるとき、実体は不動なままに
 留まり、また、実体的形相が変化する場合は、質料が不動なままに留まるからである。
 さらに、動きうる諸事物にも、不動の状態が属する。たとえば、ソクラテスは、常に座っ
 ているわけではないが、座っているときはどんな時でも、一つの場所に留まっているこ
 とは、不動に真であるように。このため、動きうるものについて、不動の知を持つこと
 を妨げるものはない。
 
 
\end{longtable}

Sententiae infra nondum translatae sunt.

\newpage


\rhead{a.2}
\begin{center}
{\Large {\bfseries ARTICULUS SECUNDUS}}\\
{\large UTRUM ANIMA PER ESSENTIAM SUAM CORPORALIA INTELLIGAT}\\
{\Large 第二項\\魂は自己の本質によって物体的なものを知性認識するか}
\end{center}

\begin{longtable}{p{21em}p{21em}}
Ad secundum sic proceditur. Videtur quod anima per essentiam suam corporalia
intelligat. Dicit enim Augustinus, X de Trin., quod anima imagines corporum
convolvit et rapit factas in semetipsa de semetipsa, dat enim eis formandis
quiddam substantiae suae. Sed per similitudines corporum corpora
intelligit. Ergo per essentiam suam, quam dat formandis talibus similitudinibus,
et de qua eas format, cognoscit corporalia.

&

第二項に対しては、以下のように進められる。魂は自らの本質によって物体的なものを認
識すると思われる。理由は以下の通り。アウグスティヌスは『三位一体論』第10巻で、
「魂は、自分自身の中で、自分自身から作られた物体の像を、結び付けたり切り離したり
する。すなわち、作られるべきそれらに自分の実体の何かを与える」と述べている。しか
るに、魂は物体の類似によって物体を知性認識する。ゆえに、作られたそのような類似に
与え、またそれからそれらの類似を形成するところの自らの本質によって、物体的なもの
を認識する。

\\[1em]


2 Praeterea, philosophus dicit, in III de anima, quod
anima quodammodo est omnia. Cum ergo simile simili cognoscatur, videtur quod
anima per seipsam corporalia cognoscat.


&

 さらに、哲学者は『デ・アニマ』第3巻で、魂はある意味で全てのものである、と述べて
 いる。ゆえに、類似したものは類似したものによって認識されるので、魂は、自分自身
 によって物体的なものを認識すると思われる。


\\[1em]



3. Praeterea, anima est superior corporalibus creaturis. Inferiora autem sunt in
superioribus eminentiori modo quam in seipsis, ut Dionysius dicit. Ergo omnes
creaturae corporeae nobiliori modo existunt in ipsa substantia animae quam in
seipsis. Per suam ergo substantiam potest creaturas corporeas cognoscere.

&



\\[1em]


Sed contra est quod Augustinus dicit, IX de Trin., quod mens corporearum rerum
notitias per sensus corporis colligit. Ipsa autem anima non est cognoscibilis
per corporis sensus. Non ergo cognoscit corporea per suam substantiam.

&



\\[1em]


Respondeo dicendum quod antiqui philosophi posuerunt quod anima per suam
essentiam cognoscit corpora. Hoc enim animis omnium communiter inditum fuit,
quod simile simili cognoscitur.

&

解答する。以下のように言われるべきである。古代の哲学者たちは、魂が自己の本質によっ
て物体的なものを認識できると考えた。理由は以下の通り。「同じものは同じものによっ
て認識される」ということが、すべての人々の心に共通に前提とされていた。


\\[1em]

Existimabant autem quod forma cogniti sit in cognoscente eo modo quo est in re
cognita. E contrario tamen Platonici posuerunt. Plato enim, quia perspexit
intellectualem animam immaterialem esse et immaterialiter cognoscere, posuit
formas rerum cognitarum immaterialiter subsistere.

&

他方で、彼らは認識されたものの形相が、認識者の中に、認識された事物においてあるあ
り方で存在すると考えていた。しかしプラトン派の人々は反対に考えた。つまりプラトン
は、知性的な魂が非質料的であり、非質料的に認識することを見て取り、認識された事物
の形相が、非質料的に自存すると考えた。

\\[1em]

Priores vero naturales, quia considerabant res cognitas esse corporeas et
materiales, posuerunt oportere res cognitas etiam in anima cognoscente
materialiter esse. Et ideo, ut animae attribuerent omnium cognitionem, posuerunt
eam habere naturam communem cum omnibus.

&

しかし、先行する自然学者たちは、認識された事物が物体的で質料的なので、認識された
事物が認識する魂の中でも質料的に存在しなければならないと考えた。ゆえに、すべての
ものの認識を玉哀史に帰するために、魂が、すべてのものに共通する本性を持つと考えた。


\\[1em]

Et quia natura principiatorum ex principiis constituitur, attribuerunt animae
naturam principii, ita quod qui dixit principium omnium esse ignem, posuit
animam esse de natura ignis; et similiter de aere et aqua. Empedocles autem, qui
posuit quatuor elementa materialia et duo moventia, ex his etiam dixit animam
esse constitutam.Et ita, cum res materialiter in anima ponerent, posuerunt omnem
cognitionem animae materialem esse, non discernentes inter intellectum et
sensum.

&

そして、根源から生じたものどもの本性は、根源から構成されるので、彼らは魂に、根源
の本性を帰した。そのようにして、万物の根源が火であるといった人は、魂が火の本性か
らなると考え、同様に、（別の人々は）空気や水からなると考えた。他方でエンペドクレ
スは、四つの質料的な元素と二つの動的原理を考え、これらから魂も構成されていると言っ
た。そして、このようにして、彼らは事物を質料的に魂の中に置いたので、魂のすべての認識
が質料的であると考え、知性と感覚を区別しなかった。
 

\\[1em]

Sed haec opinio improbatur. Primo quidem, quia in materiali principio, de quo
loquebantur, non existunt principiata nisi in potentia.

&

しかしこの意見は論駁される。第一に、彼らが語っていた質料的な原理において、原理によっ
て生じるものは、可能態においてしか存在しないからである。

\\[1em]

Non autem cognoscitur aliquid secundum quod est in potentia, sed solum secundum
quod est actu, ut patet in IX Metaphys., unde nec ipsa potentia cognoscitur nisi
per actum. Sic igitur non sufficeret attribuere animae principiorum naturam ad
hoc quod omnia cognosceret, nisi inessent ei naturae et formae singulorum
effectuum, puta ossis et carnis et aliorum huiusmodi; ut Aristoteles contra
Empedoclem argumentatur in I de anima.

&

しかし、『形而上学』第9巻で明らかなように、なにものも可能態においてある限りでは
認識されず、現実態においてある限りでのみ認識される。このことから、可能態（能力）
もまた、現実態（作用）によってしか認識されない。ゆえにこのようにして、ちょうどア
リストテレスが『デ・アニマ』第1巻でエンペドクレスに反対して論じているように、た
とえば骨や肉や、その他そのようなものの、個的な結果の本性や形相がそれに内在してい
なかったならば、魂に諸根源の本性を帰することは、万物を認識することに十分でなかっ
ただろう。

\\[1em]

Secundo quia, si oporteret rem cognitam materialiter in cognoscente existere,
nulla ratio esset quare res quae materialiter extra animam subsistunt,
cognitione carerent, puta, si anima igne cognoscit ignem, et ignis etiam qui est
extra animam, ignem cognosceret. Relinquitur ergo quod oportet materialia
cognita in cognoscente existere non materialiter, sed magis immaterialiter.

&

第二に、もし、認識された事物が質料的に認識者の中に存在しなければならなかったなら
ば、魂の外に質料的に自存する事物が認識を欠く理由がまったくなかったであろう。たと
えば、もし、魂が火によって火を認識するのであれば、魂の外にある火もまた火を認識し
たであろう。ゆえに、認識された質料的事物は認識者の中に質料的にでなく、むしろ非質
料的に存在することが残される。\footnote{認識主体と認識の媒介が同一視されているよ
うに見える。各々の人の魂が、魂の中の火によって火を認識することから、魂の外にある
火が、火によって火を認識するとはどういうことか。}

\\[1em]

Et huius ratio est, quia actus cognitionis se extendit ad ea quae sunt extra
cognoscentem, cognoscimus enim etiam ea quae extra nos sunt. Per materiam autem
determinatur forma rei ad aliquid unum. Unde manifestum est quod ratio
cognitionis ex opposito se habet ad rationem materialitatis.

&

そしてこの理由は以下の通りである。すなわち、認識の作用は、認識者の外にあるものに
まで広がる。つまり、私たちは、私たちの外にあるものを認識する。しかるに、質料によっ
て、事物の形相はある一つのものに限定される。したがって、認識の性格が、質料性の性
格の正反対にあることが明らかである。

\\[1em]

Et ideo quae non recipiunt formas nisi materialiter, nullo modo sunt
cognoscitiva, sicut plantae; ut dicitur in II libro de anima. Quanto autem
aliquid immaterialius habet formam rei cognitae, tanto perfectius
cognoscit. Unde et intellectus, qui abstrahit speciem non solum a materia, sed
etiam a materialibus conditionibus individuantibus, perfectius cognoscit quam
sensus, qui accipit formam rei cognitae sine materia quidem, sed cum
materialibus conditionibus.

&

ゆえに、質料的にしか形相を受け取らないものは、植物のように、決して認識しうるもの
ではない。これは『デ・アニマ』第2巻で言われているとおりである。しかし、何かが認
識される事物の形相をより非質料的に持てばもつほど、より完全に認識する。したがって
知性も、質料からだけでなく、個体化する質料的諸条件からもまた抽象するのだが、感覚
よりも完全に認識する。感覚は、認識された事物の形相を、質料なしに受け取るが、しか
し質料的な諸条件を伴って受け取るからである。
 

\\[1em]

Et inter ipsos sensus, visus est magis cognoscitivus, quia est minus materialis,
ut supra dictum est. Et inter ipsos intellectus, tanto quilibet est perfectior,
quanto immaterialior. Ex his ergo patet quod, si aliquis intellectus est qui per
essentiam suam cognoscit omnia, oportet quod essentia eius habeat in se
immaterialiter omnia; sicut antiqui posuerunt essentiam animae actu componi ex
principiis omnium materialium, ut cognosceret omnia.

&

そしてこれらの感覚の中で、視覚はより認識する力が強い。なぜなら、前に述べられたと
おり、質料的である度合いがより少ないからである。そして、知性の中では、より非質料
的なものの方が、より完全である。ゆえに、以上のことから、もし、ある知性が、自己の
本質を通して万物を認識するのであれば、それの本質は、自らの中に、非質料的に万物を
持っていなければならない。ちょうど、古代の人々が、魂が万物を認識するためには、魂
の本質が、現実に、すべての質料的事物の諸根源から複合されていると考えたように。
 
\\[1em]

Hoc autem est proprium Dei, ut sua essentia sit immaterialiter comprehensiva omnium, prout effectus virtute praeexistunt in causa. Solus igitur Deus per essentiam suam omnia intelligit; non autem anima humana, neque etiam Angelus.

&

 しかしこのこと、すなわち、自己の本質が、結果が原因のちからの中に先在するように
 して、非質料的に万物を包含するということは、神に固有である。ゆえに神のみが自己
 の本質によって万物を知性認識するのであり、人間の魂や、また天使もまた、そうでは
 ない。


\\[1em]

Ad primum ergo dicendum quod Augustinus ibi loquitur de visione imaginaria, quae
fit per imagines corporum. Quibus imaginibus formandis dat anima aliquid suae
substantiae, sicut subiectum datur ut informetur per aliquam formam. Et sic de
seipsa facit huiusmodi imagines, non quod anima vel aliquid animae convertatur,
ut sit haec vel illa imago; sed sicut dicitur de corpore fieri aliquid
coloratum, prout informatur colore. Et hic sensus apparet ex his quae
sequuntur. Dicit enim quod servat aliquid, scilicet non formatum tali imagine,
quod libere de specie talium imaginum iudicet, et hoc dicit esse mentem vel
intellectum. Partem autem quae informatur huiusmodi imaginibus scilicet
imaginativam, dicit esse communem nobis et bestiis.

&

第一に対しては、以下のように言われるべきである。アウグスティヌスはここで、物体の
イメージによって作られる視覚的イメージについて語っている。魂は、作られるべきそれ
らのイメージに、自分の実体の何かを与える。ちょうど、ある形相によって形相付けられ
るように、基体が与えられるようにである。このようにして、魂は自分自身からそのよう
なイメージを作るが、魂や魂の何かが、これやあのイメージであるように変わるわけでは
なく、ちょうど色について、色によって形相付けられるかぎりで、色づいた何かになると
言われるようにである。そしてこの意味は、後続することから明らかである。すなわち、
彼は、「何かを保持する」、すなわち、そのようなイメージによって形相付けられるので
はなく、「自由にそのようなイメージの種について判断する」と言う。そしてこれが、精
神ないし知性だという。ところで、このようなイメージによって形相付けられる部分、す
なわち想像的部分は、「私たちと動物に共通である」と言われる。

 

\\[1em]


[31963] Iª q. 84 a. 2 ad 2
Ad secundum dicendum quod Aristoteles non posuit animam esse actu compositam ex omnibus, sicut antiqui naturales; sed dixit quodammodo animam esse omnia, inquantum est in potentia ad omnia; per sensum quidem ad sensibilia, per intellectum vero ad intelligibilia.

&



\\[1em]


[31964] Iª q. 84 a. 2 ad 3
Ad tertium dicendum quod quaelibet creatura habet esse finitum et determinatum. Unde essentia superioris creaturae, etsi habeat quandam similitudinem inferioris creaturae prout communicant in aliquo genere, non tamen complete habet similitudinem illius, quia determinatur ad aliquam speciem, praeter quam est species inferioris creaturae. Sed essentia Dei est perfecta similitudo omnium quantum ad omnia quae in rebus inveniuntur, sicut universale principium omnium.

&



\end{longtable}
\newpage






\rhead{a.~3}
\begin{center}
{\Large {\bfseries ARTICULUS TERTIUS}}\\
{\large UTRUM ANIMA INTELLIGAT OMNIA PER SPECIES SIBI NATURALITER INDITAS}\\
{\Large 第三項\\魂は自らに本性的に与えられた形象によって万物を知性認識するか}
\end{center}

\begin{longtable}{p{21em}p{21em}}
ARTICULUS 3
[31965] Iª q. 84 a. 3 arg. 1
Ad tertium sic proceditur. Videtur quod anima intelligat omnia per species sibi naturaliter inditas. Dicit enim Gregorius, in homilia ascensionis quod homo habet commune cum Angelis intelligere. Sed Angeli intelligunt omnia per formas naturaliter inditas, unde in libro de causis dicitur quod omnis intelligentia est plena formis. Ergo et anima habet species rerum naturaliter inditas, quibus corporalia intelligit.


 
&

\\[1em]



 
[31966] Iª q. 84 a. 3 arg. 2
 Praeterea, anima intellectiva est nobilior quam materia prima corporalis. Sed materia prima est creata a Deo sub formis ad quas est in potentia. Ergo multo magis anima intellectiva est creata a Deo sub speciebus intelligibilibus. Et sic anima intelligit corporalia per species sibi naturaliter inditas.

 
 
&

\\[1em]




[31967] Iª q. 84 a. 3 arg. 3
 Praeterea, nullus potest verum respondere nisi de eo quod scit. Sed aliquis etiam idiota, non habens scientiam acquisitam, respondet verum de singulis, si tamen ordinate interrogetur, ut narratur in Menone Platonis de quodam. Ergo antequam aliquis acquirat scientiam, habet rerum cognitionem. Quod non esset nisi anima haberet species naturaliter inditas. Intelligit igitur anima res corporeas per species naturaliter inditas.

 
 
&

\\[1em]




[31968] Iª q. 84 a. 3 s. c.
 Sed contra est quod philosophus dicit, in III de anima, de intellectu loquens, quod est sicut tabula in qua nihil est scriptum.

 
 
&

\\[1em]

 Respondeo dicendum quod, cum forma sit principium actionis, oportet ut eo modo
 se habeat aliquid ad formam quae est actionis principium, quo se habet ad
 actionem illam, sicut si moveri sursum est ex levitate, oportet quod in
 potentia tantum sursum fertur, esse leve solum in potentia, quod autem actu
 sursum fertur, esse leve in actu.

 
 &

 解答する。以下のように言われるべきである。形相は作用の根源なので、あるものは、
 作用に関係するのと同じように、その作用の根源である形相に関係しなければならない。
 たとえば、もし上へ動くことが軽さに基づくならば、可能態においてしか上へ動かない
 ものは、可能態においてのみ軽いのでなければならず、また、現実に上へ動くものは、
 現実において軽いのでなければならないように。

\\[1em]

 Videmus autem quod homo est quandoque cognoscens in potentia tantum, tam
 secundum sensum quam secundum intellectum. Et de tali potentia in actum
 reducitur, ut sentiat quidem, per actiones sensibilium in sensum; ut intelligat
 autem, per disciplinam aut inventionem. Unde oportet dicere quod anima
 cognoscitiva sit in potentia tam ad similitudines quae sunt principia
 sentiendi, quam ad similitudines quae sunt principia intelligendi.
 
 &

 しかるに、私たちは、人間が、感覚においても知性においても、時として、ただ可能態
 において認識するものであることを見る。そしてそのような可能態から、感覚するため
 に、現実態へ引き出されるのは、感覚へ至る、可感的なものどもについての作用によっ
 てであり、他方で、知性認識するためには、学習と発見による。このことから、認識し
 うる魂は、感覚の根源である類似だけでなく、知性認識の根源である類似に対しても、
 可能態にあると言わなければならない。
 
\\[1em]

 Et propter hoc Aristoteles posuit quod intellectus, quo anima intelligit, non
 habet aliquas species naturaliter inditas, sed est in principio in potentia ad
 huiusmodi species omnes.

 &

 そしてアリストテレスはこのために、魂がそれによって知性認識するところの知性が、
 本性的に与えられた何らの形象も持たず、そのようなすべての形象に対して、始まりに
 おいては可能態にあると考えた。

\\[1em]

 Sed quia id quod habet actu formam, interdum non potest agere secundum formam
 propter aliquod impedimentum, sicut leve si impediatur sursum ferri; propter
 hoc Plato posuit quod intellectus hominis naturaliter est plenus omnibus
 speciebus intelligibilibus, sed per unionem corporis impeditur ne possit in
 actum exire.
 
 &

 しかし、現実に形相をもつものが、何らかの妨げによって、その形相に即して作用する
 ことができないことがある。たとえば、軽いものが、上昇することを妨げられる場合の
 ように。だから、このためにプラトンは、人間の知性が自然本性的にすべての可知的形
 象に満ちているが、身体との合一によって、現実態へ出ていくことが妨げられていると
 考えた。

\\[1em]


Sed hoc non videtur convenienter dictum.  Primo quidem quia, si habet anima
 naturalem notitiam omnium, non videtur esse possibile quod huius naturalis
 notitiae tantam oblivionem capiat, quod nesciat se huiusmodi scientiam habere,
 nullus enim homo obliviscitur ea quae naturaliter cognoscit, sicut quod omne
 totum sit maius sua parte, et alia huiusmodi. 
 
 &

 しかしこれは適切に語られているとは思えない。第一に、もし魂が万物についての自然
 本性的な知をもっていたら、自分がそのような知をもっていることを知らないほどまで
 に大きな、そのような自然本性的な知はの忘却を持つことはありえないと思われるから
 である。というのも、すべて全体はその部分よりも大きいとか、そのような、自然本性
 的に認識することを忘れる人はいないからである。

\\[1em]


 Praecipue autem hoc videtur inconveniens, si ponatur esse animae naturale
 corpori uniri, ut supra habitum est, inconveniens enim est quod naturalis
 operatio alicuius rei totaliter impediatur per id quod est sibi secundum
 naturam.

 &

 またとくに、上述のように、もし身体に合一されることが魂にとって自然本性的だとさ
 れるならば、これは不都合だと思われる。なぜなら、ある事物の自然本性的な働きが、
 自然本性に即してそれに属することによって全面的に妨げられるのは不都合だからであ
 る。


 \\[1em]


 Secundo, manifeste apparet huius positionis falsitas ex hoc quod, deficiente
 aliquo sensu, deficit scientia eorum, quae apprehenduntur secundum illum
 sensum; sicut caecus natus nullam potest habere notitiam de coloribus. Quod non
 esset, si animae essent naturaliter inditae omnium intelligibilium rationes. Et
 ideo dicendum est quod anima non cognoscit corporalia per species naturaliter
 inditas.

 
 &

 第二に、この立場の誤りは以下のことから明らかである。すなわち、何らかの感覚がな
 くなれば、その感覚に即して捉えられるものどもの知がなくなる。たとえば、生まれつ
 き目が見えない人は、色についての知をまったくもつことができない。このようなこと
 は、もし魂に、自然本性的にすべての可知的なものどもの概念が与えられていたらなかっ
 たであろう。ゆえに、魂は、自然本性的に与えられた形象で物体的なものを認識するの
 ではないと言われるべきである。

\\[1em]




[31970] Iª q. 84 a. 3 ad 1
 Ad primum ergo dicendum quod homo quidem convenit cum Angelis in intelligendo, deficit tamen ab eminentia intellectus eorum, sicut et corpora inferiora, quae tantum existunt secundum Gregorium, deficiunt ab existentia superiorum corporum. Nam materia inferiorum corporum non est completa totaliter per formam, sed est in potentia ad formas quas non habet, materia autem caelestium corporum est totaliter completa per formam, ita quod non est in potentia ad aliam formam, ut supra habitum est. Et similiter intellectus Angeli est perfectus per species intelligibiles secundum suam naturam, intellectus autem humanus est in potentia ad huiusmodi species.

 
 
&

\\[1em]




[31971] Iª q. 84 a. 3 ad 2
 Ad secundum dicendum quod materia prima habet esse substantiale per formam, et ideo oportuit quod crearetur sub aliqua forma, alioquin non esset in actu. Sub una tamen forma existens, est in potentia ad alias. Intellectus autem non habet esse substantiale per speciem intelligibilem; et ideo non est simile.

 
 
&

\\[1em]




[31972] Iª q. 84 a. 3 ad 3
 Ad tertium dicendum quod ordinata interrogatio procedit ex principiis communibus per se notis, ad propria. Per talem autem processum scientia causatur in anima addiscentis. Unde cum verum respondet de his de quibus secundo interrogatur, hoc non est quia prius ea noverit; sed quia tunc ea de novo addiscit. Nihil enim refert utrum ille qui docet, proponendo vel interrogando procedat de principiis communibus ad conclusiones, utrobique enim animus audientis certificatur de posterioribus per priora.

 
 
&


\end{longtable}
\newpage
\rhead{a.4}
\begin{center}
{\Large {\bfseries ARTICULUS 4}}\\
{\large UTRUM SPECIES INTELLIGIBILES EFFLUANT IN ANIMAM AB ALIQUIBUS FORMIS SEPARATIS}\\
{\footnotesize {\itshape De Verit.}, q.10, a.6; q.11, a.1; Qu. {\itshape de Anima}, art.15.}\\
{\Large 第四項\\知性的形象は何らかの分離した形相から魂へ流入するか}
\end{center}

\begin{longtable}{p{21em}p{21em}}

 Ad quartum sic proceditur. Videtur quod species intelligibiles effluant in animam ab aliquibus formis separatis. Omne enim quod per participationem est tale, causatur ab eo quod est per essentiam tale; sicut quod est ignitum reducitur sicut in causam in ignem. Sed anima intellectiva, secundum quod est actu intelligens, participat ipsa intelligibilia, intellectus enim in actu, quodammodo est intellectum in actu. Ergo ea quae secundum se et per essentiam suam sunt intellecta in actu, sunt causae animae intellectivae quod actu intelligat. Intellecta autem in actu per essentiam suam, sunt formae sine materia existentes. Species igitur intelligibiles quibus anima intelligit, causantur a formis aliquibus separatis.

 
&



\\[1em]




 Praeterea, intelligibilia se habent ad intellectum, sicut sensibilia ad sensum. Sed sensibilia quae sunt in actu extra animam, sunt causae specierum sensibilium quae sunt in sensu, quibus sentimus. Ergo species intelligibiles quibus intellectus noster intelligit, causantur ab aliquibus actu intelligibilibus extra animam existentibus. Huiusmodi autem non sunt nisi formae a materia separatae. Formae igitur intelligibiles intellectus nostri effluunt ab aliquibus substantiis separatis.

  
&



\\[1em]





 Praeterea, omne quod est in potentia, reducitur in actum per id quod est actu. Si ergo intellectus noster, prius in potentia existens, postmodum actu intelligat, oportet quod hoc causetur ab aliquo intellectu qui semper est in actu. Hic autem est intellectus separatus. Ergo ab aliquibus substantiis separatis causantur species intelligibiles quibus actu intelligimus.

  
&



\\[1em]





 Sed contra est quia secundum hoc sensibus non indigeremus ad intelligendum. Quod patet esse falsum ex hoc praecipue quod qui caret uno sensu, nullo modo potest habere scientiam de sensibilibus illius sensus.

  
&



\\[1em]


 Respondeo dicendum quod quidam posuerunt species intelligibiles nostri
 intellectus procedere ab aliquibus formis vel substantiis separatis. Et hoc
 dupliciter.
 
&

解答する。以下のように言われるべきである。ある人々は、私たちの知性がもつ知性的形
象が、なんらかの形相あるいは分離実体から発出すると考えた。そしてこれは二通りある。

\\[1em]

 Plato enim, sicut dictum est, posuit formas rerum sensibilium per se sine
 materia subsistentes; sicut formam hominis, quam nominabat per se hominem, et
 formam vel ideam equi, quam nominabat per se equum, et sic de aliis.

 &

 プラトンは、すでに述べられたとおり、感覚的諸事物の形相が、質料なしに自体的に自
 存すると考えた。たとえば、人間の形相は、彼はこれを自体的な人間と名付けたし、馬
 の形相ないしイデアは、これを自体的な馬と名付け、その他についても同様であった。

\\[1em]

 Has ergo formas separatas ponebat participari et ab anima nostra, et a materia
 corporali; ab anima quidem nostra ad cognoscendum, a materia vero corporali ad
 essendum; ut sicut materia corporalis per hoc quod participat ideam lapidis,
 fit hic lapis, ita intellectus noster per hoc quod participat ideam lapidis,
 fit intelligens lapidem.
  
&

ゆえに、これらの分離形相が、私たちの魂と物体的質料に分有され、そして私たちの魂に
よっては認識するために、物体的質料によっては存在するために分有されると考えた。た
とえば、ちょうど、物体的質料が、石のイデアを分有することによってこの石になるよう
に、私たちの知性は、石のイデアを分有することによって、石を認識するものになる、と
いうように。

\\[1em]


Participatio autem ideae fit per aliquam similitudinem ipsius ideae in
participante ipsam, per modum quo exemplar participatur ab exemplato. Sicut
igitur ponebat formas sensibiles quae sunt in materia corporali, effluere ab
ideis sicut quasdam earum similitudines; ita ponebat species intelligibiles
nostri intellectus esse similitudines quasdam idearum ab eis effluentes. Et
propter hoc, ut supra dictum est, scientias et definitiones ad ideas referebat.
 
&

 ところで、イデアの分有は、範型が範型から生じるものに分有されるしかたで、イデア
 自身の、イデアを分有するものにおける何らかの類似範型によって生じる。ゆえに、ちょ
 うど、物体的質料においてある感覚的形相が、イデアから、イデアの何らかの類似とし
 て流出すると考えたように、私たちの知性が持つ知性的形象は、イデアから流出した、
 何らかのイデアの類似だと考えた。このために、上述のように、学知と定義はイデアへ
 と関係づけられた。


\\[1em]

 Sed quia contra rationem rerum sensibilium est quod earum formae subsistant
 absque materiis, ut Aristoteles multipliciter probat; ideo Avicenna, hac
 positione remota, posuit omnium rerum sensibilium intelligibiles species, non
 quidem per se subsistere absque materia, sed praeexistere immaterialiter in
 intellectibus separatis; a quorum primo derivantur huiusmodi species in
 sequentem, et sic de aliis usque ad ultimum intellectum separatum, quem nominat
 intellectum agentem; a quo, ut ipse dicit, effluunt species intelligibiles in
 animas nostras, et formae sensibiles in materiam corporalem.
 
 &
 
 しかし、アリストテレスが多くのしかたで証明しているように、質料なしにそれらの形
 相が自存するということは、感覚的諸事物の性格に反するので、アヴィセンナは、この
 立場から離れて、全ての感覚的事物の知性的形象は、質料なしにそれ自体で自存するの
 ではなく、分離した知性の中に非質料的に先在すると考えた。そしてそれらのうちの第
 一のものからそのような形象が後続するものへと派生し、その他のものについても同様
 にして、最終の分離知性、彼はそれを能動知性と読んだが、にまで至るとした。そして
 彼が言うには、その能動知性から、私たちの魂へと知性的形象が、物体的質料へと感覚
 的形相が流出するとした。
 



\\[1em]

 Et sic in hoc Avicenna cum Platone concordat, quod species intelligibiles
 nostri intellectus effluunt a quibusdam formis separatis, quas tamen Plato
 dicit per se subsistere, Avicenna vero ponit eas in intelligentia
 agente. Differunt etiam quantum ad hoc, quod Avicenna ponit species
 intelligibiles non remanere in intellectu nostro postquam desinit actu
 intelligere; sed indiget ut iterato\footnote{from itero, adv, again, once
 more} se convertat ad recipiendum de novo. Unde non ponit scientiam animae
 naturaliter inditam, sicut Plato, qui ponit participationes idearum immobiliter
 in anima permanere.

 &

 このように、アヴィセンナは、私たちの知性の知性的形象が、何らかの分離形相から流
 出する点ではプラトンと一致するが、しかしプラトンはそれらが自存すると言い、アヴィ
 センナはそれらが能動知性の中にあるとする。さらに、アヴィセンナは知性的形象が、
 現実的な知性認識を止めた後に、私たちの知性の中に残らず、新たに受け取るために、
 それに向き直ることを必要とする、という点でも彼らは異なる。したがって、イデアの
 分有は魂の中に不動に留まるとするプラトンのように、知が魂に自然本性的に植え付け
 られているとは考えない。

\\[1em]

 Sed secundum hanc positionem sufficiens ratio assignari non posset quare anima
 nostra corpori uniretur. Non enim potest dici quod anima intellectiva corpori
 uniatur propter corpus, quia nec forma est propter materiam, nec motor propter
 mobile, sed potius e converso. Maxime autem videtur corpus esse necessarium
 animae intellectivae ad eius propriam operationem, quae est intelligere, quia
 secundum esse suum a corpore non dependet. Si autem anima species
 intelligibiles secundum suam naturam apta nata esset recipere per influentiam
 aliquorum separatorum principiorum tantum, et non acciperet eas ex sensibus,
 non indigeret corpore ad intelligendum, unde frustra corpori uniretur.
 
 &

 しかし、この考えによれば、魂がなぜ身体に合一されているかの十分な理由が指定され
 ることができない。すなわち、知性的な魂が身体のために身体に合一されると言われる
 ことができない。なぜなら、形相は質料のためにあるのでなく、動者も可動的なものの
 ためにあるのでもなく、むしろ逆だからである。しかるに、身体は、知性的魂に、その
 固有の働きである知性認識のために必要であると強く思われる。なぜなら、自らの存在
 に即しては、身体に依存しないからである。しかるに、もし魂が知性的形象を、自らの
 本性に即して、他の分離した諸根源の流入によってのみ受け取り、それらを感覚から受
 け取るのでないとしたら、知性認識のために身体を必要とせず、したがって、身体に無
 駄に合一されていることになっただろう。

\\[1em]

 Si autem dicatur quod indiget anima nostra sensibus ad intelligendum, quibus
 quodammodo excitetur ad consideranda ea quorum species intelligibiles a
 principiis separatis recipit; hoc non videtur sufficere. Quia huiusmodi
 excitatio non videtur necessaria animae nisi inquantum est
 consopita\footnote{con-sōpĭo, no perf., ītum, 4, v. a., to bring into an
 unconscious state, to put fast asleep, lull to sleep, to stupefy}, secundum
 Platonicos, quodammodo et obliviosa propter unionem ad corpus, et sic sensus
 non proficerent animae intellectivae nisi ad tollendum impedimentum quod animae
 provenit ex corporis unione. Remanet igitur quaerendum quae sit causa unionis
 animae ad corpus.
 
 &

 
 


 
\\[1em]

 Si autem dicatur, secundum Avicennam, quod sensus sunt animae necessarii, quia
 per eos excitatur ut convertat se ad intelligentiam agentem, a qua recipit
 species; hoc quidem non sufficit. Quia si in natura animae est ut intelligat
 per species ab intelligentia agente effluxas, sequeretur quod quandoque anima
 possit se convertere ad intelligentiam agentem ex inclinatione suae naturae,
 vel etiam excitata per alium sensum, ut convertat se ad intelligentiam agentem
 ad recipiendum species sensibilium quorum sensum aliquis non habet. Et sic
 caecus natus posset habere scientiam de coloribus, quod est manifeste
 falsum. Unde dicendum est quod species intelligibiles quibus anima nostra
 intelligit, non effluunt a formis separatis.
  
&

\\[1em]


 Ad primum ergo dicendum quod species intelligibiles quas participat noster intellectus, reducuntur sicut in primam causam in aliquod principium per suam essentiam intelligibile, scilicet in Deum. Sed ab illo principio procedunt mediantibus formis rerum sensibilium et materialium, a quibus scientiam colligimus, ut Dionysius dicit.

  
&



\\[1em]




[31979] Iª q. 84 a. 4 ad 2
 Ad secundum dicendum quod res materiales, secundum esse quod habent extra animam, possunt esse sensibiles actu; non autem actu intelligibiles. Unde non est simile de sensu et intellectu.

  
&



\\[1em]





 Ad tertium dicendum quod intellectus noster possibilis reducitur de potentia ad actum per aliquod ens actu, idest per intellectum agentem, qui est virtus quaedam animae nostrae, ut dictum est, non autem per aliquem intellectum separatum, sicut per causam proximam; sed forte sicut per causam remotam.
 
&


\end{longtable}
\newpage

\end{document}



%\rhead{a.~}
%\begin{center}
% {\Large {\bfseries }}\\
% {\large }\\
% {\footnotesize }\\
% {\Large \\}
%\end{center}
%
%\begin{longtable}{p{21em}p{21em}}
%
%&
%
% 
%
%\\
%\end{longtable}
%\newpage

ARTICULUS 5
[31981] Iª q. 84 a. 5 arg. 1
Ad quintum sic proceditur. Videtur quod anima intellectiva non cognoscat res materiales in rationibus aeternis. Id enim in quo aliquid cognoscitur, ipsum magis et per prius cognoscitur. Sed anima intellectiva hominis, in statu praesentis vitae, non cognoscit rationes aeternas, quia non cognoscit ipsum Deum, in quo rationes aeternae existunt, sed ei sicut ignoto coniungitur, ut Dionysius dicit in I cap. mysticae theologiae. Ergo anima non cognoscit omnia in rationibus aeternis.

[31982] Iª q. 84 a. 5 arg. 2
Praeterea, Rom. I, dicitur quod invisibilia Dei per ea quae facta sunt, conspiciuntur. Sed inter invisibilia Dei numerantur rationes aeternae. Ergo rationes aeternae per creaturas materiales cognoscuntur, et non e converso.

[31983] Iª q. 84 a. 5 arg. 3
Praeterea, rationes aeternae nihil aliud sunt quam ideae, dicit enim Augustinus, in libro octoginta trium quaest., quod ideae sunt rationes stabiles rerum in mente divina existentes. Si ergo dicatur quod anima intellectiva cognoscit omnia in rationibus aeternis, redibit opinio Platonis, qui posuit omnem scientiam ab ideis derivari.

[31984] Iª q. 84 a. 5 s. c.
Sed contra est quod dicit Augustinus, XII Confess., si ambo videmus verum esse quod dicis, et ambo videmus verum esse quod dico, ubi quaeso id videmus? Nec ego utique in te, nec tu in me sed ambo in ipsa, quae supra mentes nostras est, incommutabili veritate. Veritas autem incommutabilis in aeternis rationibus continetur. Ergo anima intellectiva omnia vera cognoscit in rationibus aeternis.

[31985] Iª q. 84 a. 5 co.
Respondeo dicendum quod, sicut Augustinus dicit in II de Doctr. Christ., philosophi qui vocantur, si qua forte vera et fidei nostrae accommoda dixerunt, ab eis tanquam ab iniustis possessoribus in usum nostrum vindicanda sunt. Habent enim doctrinae gentilium quaedam simulata et superstitiosa figmenta, quae unusquisque nostrum de societate gentilium exiens, debet evitare. Et ideo Augustinus, qui doctrinis Platonicorum imbutus fuerat, si qua invenit fidei accommoda in eorum dictis, assumpsit; quae vero invenit fidei nostrae adversa, in melius commutavit. Posuit autem Plato, sicut supra dictum est, formas rerum per se subsistere a materia separatas, quas ideas vocabat, per quarum participationem dicebat intellectum nostrum omnia cognoscere; ut sicut materia corporalis per participationem ideae lapidis fit lapis, ita intellectus noster per participationem eiusdem ideae cognosceret lapidem. Sed quia videtur esse alienum a fide quod formae rerum extra res per se subsistant absque materia, sicut Platonici posuerunt, dicentes per se vitam aut per se sapientiam esse quasdam substantias creatrices, ut Dionysius dicit XI cap. de Div. Nom.; ideo Augustinus, in libro octoginta trium quaest., posuit loco harum idearum quas Plato ponebat, rationes omnium creaturarum in mente divina existere, secundum quas omnia formantur, et secundum quas etiam anima humana omnia cognoscit. Cum ergo quaeritur utrum anima humana in rationibus aeternis omnia cognoscat, dicendum est quod aliquid in aliquo dicitur cognosci dupliciter. Uno modo, sicut in obiecto cognito; sicut aliquis videt in speculo ea quorum imagines in speculo resultant. Et hoc modo anima, in statu praesentis vitae, non potest videre omnia in rationibus aeternis; sed sic in rationibus aeternis cognoscunt omnia beati, qui Deum vident et omnia in ipso. Alio modo dicitur aliquid cognosci in aliquo sicut in cognitionis principio; sicut si dicamus quod in sole videntur ea quae videntur per solem. Et sic necesse est dicere quod anima humana omnia cognoscat in rationibus aeternis, per quarum participationem omnia cognoscimus. Ipsum enim lumen intellectuale quod est in nobis, nihil est aliud quam quaedam participata similitudo luminis increati, in quo continentur rationes aeternae. Unde in Psalmo IV, dicitur, multi dicunt, quis ostendit nobis bona? Cui quaestioni Psalmista respondet, dicens, signatum est super nos lumen vultus tui, domine. Quasi dicat, per ipsam sigillationem divini luminis in nobis, omnia nobis demonstrantur. Quia tamen praeter lumen intellectuale in nobis, exiguntur species intelligibiles a rebus acceptae, ad scientiam de rebus materialibus habendam; ideo non per solam participationem rationum aeternarum de rebus materialibus notitiam habemus, sicut Platonici posuerunt quod sola idearum participatio sufficit ad scientiam habendam. Unde Augustinus dicit, in IV de Trin., numquid quia philosophi documentis certissimis persuadent aeternis rationibus omnia temporalia fieri, propterea potuerunt in ipsis rationibus perspicere, vel ex ipsis colligere quot sint animalium genera, quae semina singulorum? Nonne ista omnia per locorum ac temporum historiam quaesierunt? Quod autem Augustinus non sic intellexerit omnia cognosci in rationibus aeternis, vel in incommutabili veritate, quasi ipsae rationes aeternae videantur, patet per hoc quod ipse dicit in libro octoginta trium quaest., quod rationalis anima non omnis et quaelibet, sed quae sancta et pura fuerit, asseritur illi visioni, scilicet rationum aeternarum, esse idonea; sicut sunt animae beatorum.

[31986] Iª q. 84 a. 5 ad arg.
Et per haec patet responsio ad obiecta.

%\rhead{a.~}
%\begin{center}
% {\Large {\bfseries }}\\
% {\large }\\
% {\footnotesize }\\
% {\Large \\}
%\end{center}
%
%\begin{longtable}{p{21em}p{21em}}
%
%&
%
% 
%
%\\
%\end{longtable}
%\newpage

ARTICULUS 6
[31987] Iª q. 84 a. 6 arg. 1
Ad sextum sic proceditur. Videtur quod intellectiva cognitio non accipiatur a rebus sensibilibus. Dicit enim Augustinus, in libro octoginta trium quaest., quod non est expectanda sinceritas veritatis a corporis sensibus. Et hoc probat dupliciter. Uno modo, per hoc quod omne quod corporeus sensus attingit, sine ulla intermissione temporis commutatur, quod autem non manet, percipi non potest. Alio modo, per hoc quod omnia quae per corpus sentimus, etiam cum non adsunt sensibus, imagines tamen eorum patimur, ut in somno vel furore; non autem sensibus discernere valemus utrum ipsa sensibilia, vel imagines eorum falsas sentiamus. Nihil autem percipi potest quod a falso non discernitur. Et sic concludit quod non est expectanda veritas a sensibus. Sed cognitio intellectualis est apprehensiva veritatis. Non ergo cognitio intellectualis est expectanda a sensibus.

[31988] Iª q. 84 a. 6 arg. 2
Praeterea, Augustinus dicit, XII super Gen. ad Litt., non est putandum facere aliquid corpus in spiritum, tanquam spiritus corpori facienti materiae vice subdatur, omni enim modo praestantior est qui facit, ea re de qua aliquid facit. Unde concludit quod imaginem corporis non corpus in spiritu, sed ipse spiritus in seipso facit. Non ergo intellectualis cognitio a sensibilibus derivatur.

[31989] Iª q. 84 a. 6 arg. 3
Praeterea, effectus non se extendit ultra virtutem suae causae. Sed intellectualis cognitio se extendit ultra sensibilia, intelligimus enim quaedam quae sensu percipi non possunt. Intellectualis ergo cognitio non derivatur a rebus sensibilibus.

[31990] Iª q. 84 a. 6 s. c.
Sed contra est quod philosophus probat, I Metaphys., et in fine Poster., quod principium nostrae cognitionis est a sensu.

[31991] Iª q. 84 a. 6 co.
Respondeo dicendum quod circa istam quaestionem triplex fuit philosophorum opinio. Democritus enim posuit quod nulla est alia causa cuiuslibet nostrae cognitionis, nisi cum ab his corporibus quae cogitamus, veniunt atque intrant imagines in animas nostras; ut Augustinus dicit in epistola sua ad Dioscorum. Et Aristoteles etiam dicit, in libro de Somn. et Vigil., quod Democritus posuit cognitionem fieri per idola et defluxiones. Et huius positionis ratio fuit, quia tam ipse Democritus quam alii antiqui naturales non ponebant intellectum differre a sensu, ut Aristoteles dicit in libro de anima. Et ideo, quia sensus immutatur a sensibili, arbitrabantur omnem nostram cognitionem fieri per solam immutationem a sensibilibus. Quam quidem immutationem Democritus asserebat fieri per imaginum defluxiones. Plato vero e contrario posuit intellectum differre a sensu; et intellectum quidem esse virtutem immaterialem organo corporeo non utentem in suo actu. Et quia incorporeum non potest immutari a corporeo, posuit quod cognitio intellectualis non fit per immutationem intellectus a sensibilibus, sed per participationem formarum intelligibilium separatarum, ut dictum est. Sensum etiam posuit virtutem quandam per se operantem. Unde nec ipse sensus, cum sit quaedam vis spiritualis, immutatur a sensibilibus, sed organa sensuum a sensibilibus immutantur, ex qua immutatione anima quodammodo excitatur ut in se species sensibilium formet. Et hanc opinionem tangere videtur Augustinus, XII super Gen. ad Litt., ubi dicit quod corpus non sentit, sed anima per corpus, quo velut nuntio utitur ad formandum in seipsa quod extrinsecus nuntiatur. Sic igitur secundum Platonis opinionem, neque intellectualis cognitio a sensibili procedit, neque etiam sensibilis totaliter a sensibilibus rebus; sed sensibilia excitant animam sensibilem ad sentiendum, et similiter sensus excitant animam intellectivam ad intelligendum. Aristoteles autem media via processit. Posuit enim cum Platone intellectum differre a sensu. Sed sensum posuit propriam operationem non habere sine communicatione corporis; ita quod sentire non sit actus animae tantum, sed coniuncti. Et similiter posuit de omnibus operationibus sensitivae partis. Quia igitur non est inconveniens quod sensibilia quae sunt extra animam, causent aliquid in coniunctum, in hoc Aristoteles cum Democrito concordavit, quod operationes sensitivae partis causentur per impressionem sensibilium in sensum, non per modum defluxionis, ut Democritus posuit, sed per quandam operationem. Nam et Democritus omnem actionem fieri posuit per influxionem atomorum, ut patet in I de Generat. Intellectum vero posuit Aristoteles habere operationem absque communicatione corporis. Nihil autem corporeum imprimere potest in rem incorpoream. Et ideo ad causandam intellectualem operationem, secundum Aristotelem, non sufficit sola impressio sensibilium corporum, sed requiritur aliquid nobilius, quia agens est honorabilius patiente, ut ipse dicit. Non tamen ita quod intellectualis operatio causetur in nobis ex sola impressione aliquarum rerum superiorum, ut Plato posuit, sed illud superius et nobilius agens quod vocat intellectum agentem, de quo iam supra diximus, facit phantasmata a sensibus accepta intelligibilia in actu, per modum abstractionis cuiusdam. Secundum hoc ergo, ex parte phantasmatum intellectualis operatio a sensu causatur. Sed quia phantasmata non sufficiunt immutare intellectum possibilem, sed oportet quod fiant intelligibilia actu per intellectum agentem; non potest dici quod sensibilis cognitio sit totalis et perfecta causa intellectualis cognitionis, sed magis quodammodo est materia causae.

[31992] Iª q. 84 a. 6 ad 1
Ad primum ergo dicendum quod per verba illa Augustini datur intelligi quod veritas non sit totaliter a sensibus expectanda. Requiritur enim lumen intellectus agentis, per quod immutabiliter veritatem in rebus mutabilibus cognoscamus, et discernamus ipsas res a similitudinibus rerum.

[31993] Iª q. 84 a. 6 ad 2
Ad secundum dicendum quod Augustinus ibi non loquitur de intellectuali cognitione, sed de imaginaria. Et quia, secundum Platonis opinionem, vis imaginaria habet operationem quae est animae solius; eadem ratione usus est Augustinus ad ostendendum quod corpora non imprimunt suas similitudines in vim imaginariam, sed hoc facit ipsa anima, qua utitur Aristoteles ad probandum intellectum agentem esse aliquid separatum, quia scilicet agens est honorabilius patiente. Et procul dubio oportet, secundum hanc positionem, in vi imaginativa ponere non solum potentiam passivam, sed etiam activam. Sed si ponamus, secundum opinionem Aristotelis, quod actio virtutis imaginativae sit coniuncti, nulla sequitur difficultas, quia corpus sensibile est nobilius organo animalis, secundum hoc quod comparatur ad ipsum ut ens in actu ad ens in potentia, sicut coloratum in actu ad pupillam, quae colorata est in potentia. Posset tamen dici quod, quamvis prima immutatio virtutis imaginariae sit per motum sensibilium, quia phantasia est motus factus secundum sensum, ut dicitur in libro de anima; tamen est quaedam operatio animae in homine quae dividendo et componendo format diversas rerum imagines, etiam quae non sunt a sensibus acceptae. Et quantum ad hoc possunt accipi verba Augustini.

[31994] Iª q. 84 a. 6 ad 3
Ad tertium dicendum quod sensitiva cognitio non est tota causa intellectualis cognitionis. Et ideo non est mirum si intellectualis cognitio ultra sensitivam se extendit.

%\rhead{a.~}
%\begin{center}
% {\Large {\bfseries }}\\
% {\large }\\
% {\footnotesize }\\
% {\Large \\}
%\end{center}
%
%\begin{longtable}{p{21em}p{21em}}
%
%&
%
% 
%
%\\
%\end{longtable}
%\newpage

ARTICULUS 7
[31995] Iª q. 84 a. 7 arg. 1
Ad septimum sic proceditur. Videtur quod intellectus possit actu intelligere per species intelligibiles quas penes se habet, non convertendo se ad phantasmata. Intellectus enim fit in actu per speciem intelligibilem qua informatur. Sed intellectum esse in actu, est ipsum intelligere. Ergo species intelligibiles sufficiunt ad hoc quod intellectus actu intelligat, absque hoc quod ad phantasmata se convertat.

[31996] Iª q. 84 a. 7 arg. 2
Praeterea, magis dependet imaginatio a sensu, quam intellectus ab imaginatione. Sed imaginatio potest imaginari actu, absentibus sensibilibus. Ergo multo magis intellectus potest intelligere actu, non convertendo se ad phantasmata.

[31997] Iª q. 84 a. 7 arg. 3
Praeterea, incorporalium non sunt aliqua phantasmata, quia imaginatio tempus et continuum non transcendit. Si ergo intellectus noster non posset aliquid intelligere in actu nisi converteretur ad phantasmata, sequeretur quod non posset intelligere incorporeum aliquid. Quod patet esse falsum, intelligimus enim veritatem ipsam, et Deum et Angelos.

[31998] Iª q. 84 a. 7 s. c.
Sed contra est quod philosophus dicit, in III de anima, quod nihil sine phantasmate intelligit anima.

[31999] Iª q. 84 a. 7 co.
Respondeo dicendum quod impossibile est intellectum nostrum, secundum praesentis vitae statum, quo passibili corpori coniungitur, aliquid intelligere in actu, nisi convertendo se ad phantasmata. Et hoc duobus indiciis apparet. Primo quidem quia, cum intellectus sit vis quaedam non utens corporali organo, nullo modo impediretur in suo actu per laesionem alicuius corporalis organi, si non requireretur ad eius actum actus alicuius potentiae utentis organo corporali. Utuntur autem organo corporali sensus et imaginatio et aliae vires pertinentes ad partem sensitivam. Unde manifestum est quod ad hoc quod intellectus actu intelligat, non solum accipiendo scientiam de novo, sed etiam utendo scientia iam acquisita, requiritur actus imaginationis et ceterarum virtutum. Videmus enim quod, impedito actu virtutis imaginativae per laesionem organi, ut in phreneticis; et similiter impedito actu memorativae virtutis, ut in lethargicis; impeditur homo ab intelligendo in actu etiam ea quorum scientiam praeaccepit. Secundo, quia hoc quilibet in seipso experiri potest, quod quando aliquis conatur aliquid intelligere, format aliqua phantasmata sibi per modum exemplorum, in quibus quasi inspiciat quod intelligere studet. Et inde est etiam quod quando alium volumus facere aliquid intelligere, proponimus ei exempla, ex quibus sibi phantasmata formare possit ad intelligendum. Huius autem ratio est, quia potentia cognoscitiva proportionatur cognoscibili. Unde intellectus angelici, qui est totaliter a corpore separatus, obiectum proprium est substantia intelligibilis a corpore separata; et per huiusmodi intelligibilia materialia cognoscit. Intellectus autem humani, qui est coniunctus corpori, proprium obiectum est quidditas sive natura in materia corporali existens; et per huiusmodi naturas visibilium rerum etiam in invisibilium rerum aliqualem cognitionem ascendit. De ratione autem huius naturae est, quod in aliquo individuo existat, quod non est absque materia corporali, sicut de ratione naturae lapidis est quod sit in hoc lapide, et de ratione naturae equi quod sit in hoc equo, et sic de aliis. Unde natura lapidis, vel cuiuscumque materialis rei, cognosci non potest complete et vere, nisi secundum quod cognoscitur ut in particulari existens. Particulare autem apprehendimus per sensum et imaginationem. Et ideo necesse est ad hoc quod intellectus actu intelligat suum obiectum proprium, quod convertat se ad phantasmata, ut speculetur naturam universalem in particulari existentem. Si autem proprium obiectum intellectus nostri esset forma separata; vel si naturae rerum sensibilium subsisterent non in particularibus, secundum Platonicos; non oporteret quod intellectus noster semper intelligendo converteret se ad phantasmata.

[32000] Iª q. 84 a. 7 ad 1
Ad primum ergo dicendum quod species conservatae in intellectu possibili, in eo existunt habitualiter quando actu non intelligit, sicut supra dictum est. Unde ad hoc quod intelligamus in actu, non sufficit ipsa conservatio specierum; sed oportet quod eis utamur secundum quod convenit rebus quarum sunt species, quae sunt naturae in particularibus existentes.

[32001] Iª q. 84 a. 7 ad 2
Ad secundum dicendum quod etiam ipsum phantasma est similitudo rei particularis, unde non indiget imaginatio aliqua alia similitudine particularis, sicut indiget intellectus.

[32002] Iª q. 84 a. 7 ad 3
Ad tertium dicendum quod incorporea, quorum non sunt phantasmata, cognoscuntur a nobis per comparationem ad corpora sensibilia, quorum sunt phantasmata. Sicut veritatem intelligimus ex consideratione rei circa quam veritatem speculamur; Deum autem, ut Dionysius dicit, cognoscimus ut causam, et per excessum, et per remotionem; alias etiam incorporeas substantias, in statu praesentis vitae, cognoscere non possumus nisi per remotionem, vel aliquam comparationem ad corporalia. Et ideo cum de huiusmodi aliquid intelligimus, necesse habemus converti ad phantasmata corporum, licet ipsorum non sint phantasmata.



\rhead{a.8}
\begin{center}
{\Large {\bfseries }}\\
{\large }\\
{\footnotesize }\\
{\Large \\}
\end{center}

\begin{longtable}{p{21em}p{21em}}

&



\\
\end{longtable}
\newpage

ARTICULUS 8
[32003] Iª q. 84 a. 8 arg. 1
Ad octavum sic proceditur. Videtur quod iudicium intellectus non impediatur per ligamentum sensus. Superius enim non dependet ab inferiori. Sed iudicium intellectus est supra sensum. Ergo iudicium intellectus non impeditur per ligamentum sensus.

[32004] Iª q. 84 a. 8 arg. 2
Praeterea, syllogizare est actus intellectus. In somno autem ligatur sensus, ut dicitur in libro de Somn. et Vig.; contingit tamen quandoque quod aliquis dormiens syllogizat. Ergo non impeditur iudicium intellectus per ligamentum sensus.

[32005] Iª q. 84 a. 8 s. c.
Sed contra est quod in dormiendo ea quae contra licitos mores contingunt, non imputantur ad peccatum; ut Augustinus in XII super Gen. ad Litt. dicit. Hoc autem non esset si homo in dormiendo liberum usum rationis et intellectus haberet. Ergo impeditur rationis usus per ligamentum sensus.

[32006] Iª q. 84 a. 8 co.
Respondeo dicendum quod, sicut dictum est, proprium obiectum intellectui nostro proportionatum est natura rei sensibilis. Iudicium autem perfectum de re aliqua dari non potest, nisi ea omnia quae ad rem pertinent cognoscantur, et praecipue si ignoretur id quod est terminus et finis iudicii. Dicit autem philosophus, in III de caelo, quod sicut finis factivae scientiae est opus, ita naturalis scientiae finis est quod videtur principaliter secundum sensum, faber enim non quaerit cognitionem cultelli nisi propter opus, ut operetur hunc particularem cultellum; et similiter naturalis non quaerit cognoscere naturam lapidis et equi, nisi ut sciat rationes eorum quae videntur secundum sensum. Manifestum est autem quod non posset esse perfectum iudicium fabri de cultello, si opus ignoraret, et similiter non potest esse perfectum iudicium scientiae naturalis de rebus naturalibus, si sensibilia ignorentur. Omnia autem quae in praesenti statu intelligimus, cognoscuntur a nobis per comparationem ad res sensibiles naturales. Unde impossibile est quod sit in nobis iudicium intellectus perfectum, cum ligamento sensus, per quem res sensibiles cognoscimus.

[32007] Iª q. 84 a. 8 ad 1
Ad primum ergo dicendum quod, quamvis intellectus sit superior sensu, accipit tamen aliquo modo a sensu, et eius obiecta prima et principalia in sensibilibus fundantur. Et ideo necesse est quod impediatur iudicium intellectus ex ligamento sensus.

[32008] Iª q. 84 a. 8 ad 2
Ad secundum dicendum quod sensus ligatur in dormientibus propter evaporationes quasdam et fumositates resolutas, ut dicitur in libro de Somn. et Vig. Et ideo secundum dispositionem huiusmodi evaporationum, contingit esse ligamentum sensus maius vel minus. Quando enim multus fuerit motus vaporum, ligatur non solum sensus, sed etiam imaginatio, ita ut nulla appareant phantasmata; sicut praecipue accidit cum aliquis incipit dormire post multum cibum et potum. Si vero motus vaporum aliquantulum fuerit remissior, apparent phantasmata, sed distorta et inordinata; sicut accidit in febricitantibus. Si vero adhuc magis motus sedetur, apparent phantasmata ordinata; sicut maxime solet contingere in fine dormitionis, et in hominibus sobriis et habentibus fortem imaginationem. Si autem motus vaporum fuerit modicus, non solum imaginatio remanet libera, sed etiam ipse sensus communis ex parte solvitur; ita quod homo iudicat interdum in dormiendo ea quae videt somnia esse, quasi diiudicans inter res et rerum similitudines. Sed tamen ex aliqua parte remanet sensus communis ligatus; et ideo, licet aliquas similitudines discernat a rebus, tamen semper in aliquibus decipitur. Sic igitur per modum quo sensus solvitur et imaginatio in dormiendo, liberatur et iudicium intellectus, non tamen ex toto. Unde illi qui dormiendo syllogizant, cum excitantur, semper recognoscunt se in aliquo defecisse.

